% --- Archon physics formalization source begin ---
% archon:physics
% archon:covers PhyXMiniProblems/problem_phyx_mini_0268.lean
% archon:source-report reports/phyx_mini/problem_phyx_mini_0268.source.json

\chapter{Physics problem phyx\_mini\_0268}
\label{ch:PhyXMiniProblems_problem_phyx_mini_0268}

\paragraph{Problem source.}
\#\# Physical scenario

In the overhead view of figure, a long uniform rod of mass \$0.600\textbackslash{} kg\$ is free to rotate in a horizontal plane about a vertical axis through its center. A spring with force constant \$k = 1850\textbackslash{} N/m\$ is connected horizontally between one end of the rod and a fixed wall. When the rod is in equilibrium, it is parallel to the wall.

\#\# Question

What is the period of the small oscillations that result when the rod is rotated slightly and released?

\#\# Answer choices

- A: A: 0.0562 s

- B: B: 0.0587 s

- C: C: 0.0653 s

- D: D: 0.0674 s

\#\# Figure caption (auxiliary; use the image as primary evidence)

The image shows a mechanical system diagram. There is a horizontal wall at the top. A spring labeled "k" is attached to the wall on one end, with the other end connected to one side of a red rod. The rod is positioned at an angle, and a black dot on the rod is labeled as the "Rotation axis." The image illustrates a rotational mechanical setup, possibly representing a torsional spring system.

\#\# Dataset metadata

Resonance and Harmonics; Physical Model Grounding Reasoning, Numerical Reasoning

\paragraph{Recorded answer/context.}
C: C: 0.0653 s

\paragraph{Figure/image path.}
/root/proposal\_for\_physic/science-mango/phyx\_mini\_run/phyx\_data/test\_image/268.png

\paragraph{Formalization target.}
create a compiling Lean file with sorry bodies at `PhyXMiniProblems/problem\_phyx\_mini\_0268.lean`. The Lean declarations must preserve the physical quantities, dimensions or dimensional roles, figure labels, governing-law hypotheses, and final relation expressed by this problem.
Use Mathlib/Physlib names found through LeanExplore where available. If a physics API is missing, introduce faithful local abstractions rather than scalar placeholder aliases.

\begin{theorem}[Physics formalization target]
\label{thm:physics:phyx_mini_0268:target}
\lean{PhyXMiniProblems.ProblemPhyXMini0268.problem_phyx_mini_0268}
\uses{def:physics:phyx-mini-0268:phyxminiproblems-problemphyxmini0268-massreadout, def:physics:phyx-mini-0268:phyxminiproblems-problemphyxmini0268-timereadout, def:physics:phyx-mini-0268:phyxminiproblems-problemphyxmini0268-springconstantreadout, def:physics:phyx-mini-0268:phyxminiproblems-problemphyxmini0268-centerpivotedrodspringsetup, def:physics:phyx-mini-0268:phyxminiproblems-problemphyxmini0268-matchesproblemandfigurereadouts, def:physics:phyx-mini-0268:phyxminiproblems-problemphyxmini0268-hasphysicalrodspringparameters, def:physics:phyx-mini-0268:phyxminiproblems-problemphyxmini0268-satisfiesuniformrodinertialaw, def:physics:phyx-mini-0268:phyxminiproblems-problemphyxmini0268-satisfiescenteredendpointgeometry, def:physics:phyx-mini-0268:phyxminiproblems-problemphyxmini0268-satisfieslinearizedrodspringlaws, def:physics:phyx-mini-0268:phyxminiproblems-problemphyxmini0268-satisfiesidealsmalloscillationlaw, def:physics:phyx-mini-0268:phyxminiproblems-problemphyxmini0268-matchesanswerchoice}
The assigned autoformalize agent should translate this physics problem into Lean declarations in the covered file, with theorem and lemma proof bodies written as `by sorry`.
\end{theorem}
\begin{proof}
This is an autoformalization task, not a proof task. Produce statements that can later be proved by the physics prover without weakening the physical model.
\end{proof}

% --- Archon declaration topology begin ---
\section{Declaration topology}
\begin{definition}[Mass Quantity]
  \label{def:physics:phyx-mini-0268:phyxminiproblems-problemphyxmini0268-massquantity}
  \lean{PhyXMiniProblems.ProblemPhyXMini0268.MassQuantity}
  A nonnegative physical mass, independent of the unit used to read it
\end{definition}
\begin{proof}
  This is the declared model definition of Mass Quantity.
\end{proof}
\begin{definition}[Length Quantity]
  \label{def:physics:phyx-mini-0268:phyxminiproblems-problemphyxmini0268-lengthquantity}
  \lean{PhyXMiniProblems.ProblemPhyXMini0268.LengthQuantity}
  A nonnegative physical length
\end{definition}
\begin{proof}
  This is the declared model definition of Length Quantity.
\end{proof}
\begin{definition}[Signed Length Quantity]
  \label{def:physics:phyx-mini-0268:phyxminiproblems-problemphyxmini0268-signedlengthquantity}
  \lean{PhyXMiniProblems.ProblemPhyXMini0268.SignedLengthQuantity}
  A signed displacement along the spring axis
\end{definition}
\begin{proof}
  This is the declared model definition of Signed Length Quantity.
\end{proof}
\begin{definition}[Time Quantity]
  \label{def:physics:phyx-mini-0268:phyxminiproblems-problemphyxmini0268-timequantity}
  \lean{PhyXMiniProblems.ProblemPhyXMini0268.TimeQuantity}
  A nonnegative physical duration
\end{definition}
\begin{proof}
  This is the declared model definition of Time Quantity.
\end{proof}
\begin{definition}[Spring Constant Quantity]
  \label{def:physics:phyx-mini-0268:phyxminiproblems-problemphyxmini0268-springconstantquantity}
  \lean{PhyXMiniProblems.ProblemPhyXMini0268.SpringConstantQuantity}
  A spring force constant has dimension force per length, or mass/time²
\end{definition}
\begin{proof}
  This is the declared model definition of Spring Constant Quantity.
\end{proof}
\begin{definition}[Signed Force Quantity]
  \label{def:physics:phyx-mini-0268:phyxminiproblems-problemphyxmini0268-signedforcequantity}
  \lean{PhyXMiniProblems.ProblemPhyXMini0268.SignedForceQuantity}
  A signed force component along the spring axis
\end{definition}
\begin{proof}
  This is the declared model definition of Signed Force Quantity.
\end{proof}
\begin{definition}[Moment Of Inertia Quantity]
  \label{def:physics:phyx-mini-0268:phyxminiproblems-problemphyxmini0268-momentofinertiaquantity}
  \lean{PhyXMiniProblems.ProblemPhyXMini0268.MomentOfInertiaQuantity}
  Moment of inertia has dimension mass times length squared
\end{definition}
\begin{proof}
  This is the declared model definition of Moment Of Inertia Quantity.
\end{proof}
\begin{definition}[Signed Torque Quantity]
  \label{def:physics:phyx-mini-0268:phyxminiproblems-problemphyxmini0268-signedtorquequantity}
  \lean{PhyXMiniProblems.ProblemPhyXMini0268.SignedTorqueQuantity}
  A signed torque about the vertical rotation axis
\end{definition}
\begin{proof}
  This is the declared model definition of Signed Torque Quantity.
\end{proof}
\begin{definition}[Restoring Coefficient Quantity]
  \label{def:physics:phyx-mini-0268:phyxminiproblems-problemphyxmini0268-restoringcoefficientquantity}
  \lean{PhyXMiniProblems.ProblemPhyXMini0268.RestoringCoefficientQuantity}
  A nonnegative linear restoring-torque coefficient per dimensionless radian
\end{definition}
\begin{proof}
  This is the declared model definition of Restoring Coefficient Quantity.
\end{proof}
\begin{definition}[Angular Acceleration Quantity]
  \label{def:physics:phyx-mini-0268:phyxminiproblems-problemphyxmini0268-angularaccelerationquantity}
  \lean{PhyXMiniProblems.ProblemPhyXMini0268.AngularAccelerationQuantity}
  Angular acceleration has inverse-time-squared dimension because radians are dimensionless
\end{definition}
\begin{proof}
  This is the declared model definition of Angular Acceleration Quantity.
\end{proof}
\begin{definition}[Angular Frequency Quantity]
  \label{def:physics:phyx-mini-0268:phyxminiproblems-problemphyxmini0268-angularfrequencyquantity}
  \lean{PhyXMiniProblems.ProblemPhyXMini0268.AngularFrequencyQuantity}
  A nonnegative angular frequency, with inverse-time dimension
\end{definition}
\begin{proof}
  This is the declared model definition of Angular Frequency Quantity.
\end{proof}
\begin{definition}[mass Readout]
  \label{def:physics:phyx-mini-0268:phyxminiproblems-problemphyxmini0268-massreadout}
  \lean{PhyXMiniProblems.ProblemPhyXMini0268.massReadout}
  \uses{def:physics:phyx-mini-0268:phyxminiproblems-problemphyxmini0268-massquantity}
  Read a physical mass in a selected mass unit
\end{definition}
\begin{proof}
  This is the declared model definition of mass Readout.
\end{proof}
\begin{definition}[length Readout]
  \label{def:physics:phyx-mini-0268:phyxminiproblems-problemphyxmini0268-lengthreadout}
  \lean{PhyXMiniProblems.ProblemPhyXMini0268.lengthReadout}
  \uses{def:physics:phyx-mini-0268:phyxminiproblems-problemphyxmini0268-lengthquantity}
  Read a physical length in a selected length unit
\end{definition}
\begin{proof}
  This is the declared model definition of length Readout.
\end{proof}
\begin{definition}[signed Length Readout]
  \label{def:physics:phyx-mini-0268:phyxminiproblems-problemphyxmini0268-signedlengthreadout}
  \lean{PhyXMiniProblems.ProblemPhyXMini0268.signedLengthReadout}
  \uses{def:physics:phyx-mini-0268:phyxminiproblems-problemphyxmini0268-signedlengthquantity}
  Read a signed displacement in a selected length unit
\end{definition}
\begin{proof}
  This is the declared model definition of signed Length Readout.
\end{proof}
\begin{definition}[time Readout]
  \label{def:physics:phyx-mini-0268:phyxminiproblems-problemphyxmini0268-timereadout}
  \lean{PhyXMiniProblems.ProblemPhyXMini0268.timeReadout}
  \uses{def:physics:phyx-mini-0268:phyxminiproblems-problemphyxmini0268-timequantity}
  Read a physical duration in a selected time unit
\end{definition}
\begin{proof}
  This is the declared model definition of time Readout.
\end{proof}
\begin{definition}[spring Constant Readout]
  \label{def:physics:phyx-mini-0268:phyxminiproblems-problemphyxmini0268-springconstantreadout}
  \lean{PhyXMiniProblems.ProblemPhyXMini0268.springConstantReadout}
  \uses{def:physics:phyx-mini-0268:phyxminiproblems-problemphyxmini0268-springconstantquantity}
  Read a spring constant in coherent selected mass and time units
\end{definition}
\begin{proof}
  This is the declared model definition of spring Constant Readout.
\end{proof}
\begin{definition}[signed Force Readout]
  \label{def:physics:phyx-mini-0268:phyxminiproblems-problemphyxmini0268-signedforcereadout}
  \lean{PhyXMiniProblems.ProblemPhyXMini0268.signedForceReadout}
  \uses{def:physics:phyx-mini-0268:phyxminiproblems-problemphyxmini0268-signedforcequantity}
  Read a signed force in coherent mechanical units
\end{definition}
\begin{proof}
  This is the declared model definition of signed Force Readout.
\end{proof}
\begin{definition}[moment Of Inertia Readout]
  \label{def:physics:phyx-mini-0268:phyxminiproblems-problemphyxmini0268-momentofinertiareadout}
  \lean{PhyXMiniProblems.ProblemPhyXMini0268.momentOfInertiaReadout}
  \uses{def:physics:phyx-mini-0268:phyxminiproblems-problemphyxmini0268-momentofinertiaquantity}
  Read a moment of inertia in coherent mass and length units
\end{definition}
\begin{proof}
  This is the declared model definition of moment Of Inertia Readout.
\end{proof}
\begin{definition}[signed Torque Readout]
  \label{def:physics:phyx-mini-0268:phyxminiproblems-problemphyxmini0268-signedtorquereadout}
  \lean{PhyXMiniProblems.ProblemPhyXMini0268.signedTorqueReadout}
  \uses{def:physics:phyx-mini-0268:phyxminiproblems-problemphyxmini0268-signedtorquequantity}
  Read a signed torque in coherent mechanical units
\end{definition}
\begin{proof}
  This is the declared model definition of signed Torque Readout.
\end{proof}
\begin{definition}[restoring Coefficient Readout]
  \label{def:physics:phyx-mini-0268:phyxminiproblems-problemphyxmini0268-restoringcoefficientreadout}
  \lean{PhyXMiniProblems.ProblemPhyXMini0268.restoringCoefficientReadout}
  \uses{def:physics:phyx-mini-0268:phyxminiproblems-problemphyxmini0268-restoringcoefficientquantity}
  Read a restoring coefficient in coherent mechanical units
\end{definition}
\begin{proof}
  This is the declared model definition of restoring Coefficient Readout.
\end{proof}
\begin{definition}[angular Acceleration Readout]
  \label{def:physics:phyx-mini-0268:phyxminiproblems-problemphyxmini0268-angularaccelerationreadout}
  \lean{PhyXMiniProblems.ProblemPhyXMini0268.angularAccelerationReadout}
  \uses{def:physics:phyx-mini-0268:phyxminiproblems-problemphyxmini0268-angularaccelerationquantity}
  Read angular acceleration in inverse units of a selected time unit
\end{definition}
\begin{proof}
  This is the declared model definition of angular Acceleration Readout.
\end{proof}
\begin{definition}[angular Frequency Readout]
  \label{def:physics:phyx-mini-0268:phyxminiproblems-problemphyxmini0268-angularfrequencyreadout}
  \lean{PhyXMiniProblems.ProblemPhyXMini0268.angularFrequencyReadout}
  \uses{def:physics:phyx-mini-0268:phyxminiproblems-problemphyxmini0268-angularfrequencyquantity}
  Read angular frequency in inverse units of a selected time unit
\end{definition}
\begin{proof}
  This is the declared model definition of angular Frequency Readout.
\end{proof}
\begin{definition}[mass In Kilograms]
  \label{def:physics:phyx-mini-0268:phyxminiproblems-problemphyxmini0268-massinkilograms}
  \lean{PhyXMiniProblems.ProblemPhyXMini0268.massInKilograms}
  \uses{def:physics:phyx-mini-0268:phyxminiproblems-problemphyxmini0268-massquantity, def:physics:phyx-mini-0268:phyxminiproblems-problemphyxmini0268-massreadout}
  SI kilogram readout of the rod mass
\end{definition}
\begin{proof}
  This is the declared model definition of mass In Kilograms.
\end{proof}
\begin{definition}[length In Meters]
  \label{def:physics:phyx-mini-0268:phyxminiproblems-problemphyxmini0268-lengthinmeters}
  \lean{PhyXMiniProblems.ProblemPhyXMini0268.lengthInMeters}
  \uses{def:physics:phyx-mini-0268:phyxminiproblems-problemphyxmini0268-lengthquantity, def:physics:phyx-mini-0268:phyxminiproblems-problemphyxmini0268-lengthreadout}
  SI meter readout of a length
\end{definition}
\begin{proof}
  This is the declared model definition of length In Meters.
\end{proof}
\begin{definition}[time In Seconds]
  \label{def:physics:phyx-mini-0268:phyxminiproblems-problemphyxmini0268-timeinseconds}
  \lean{PhyXMiniProblems.ProblemPhyXMini0268.timeInSeconds}
  \uses{def:physics:phyx-mini-0268:phyxminiproblems-problemphyxmini0268-timequantity, def:physics:phyx-mini-0268:phyxminiproblems-problemphyxmini0268-timereadout}
  SI second readout of a duration
\end{definition}
\begin{proof}
  This is the declared model definition of time In Seconds.
\end{proof}
\begin{definition}[spring Constant In Newtons Per Meter]
  \label{def:physics:phyx-mini-0268:phyxminiproblems-problemphyxmini0268-springconstantinnewtonspermeter}
  \lean{PhyXMiniProblems.ProblemPhyXMini0268.springConstantInNewtonsPerMeter}
  \uses{def:physics:phyx-mini-0268:phyxminiproblems-problemphyxmini0268-springconstantquantity, def:physics:phyx-mini-0268:phyxminiproblems-problemphyxmini0268-springconstantreadout}
  SI spring-constant readout, numerically in newtons per meter
\end{definition}
\begin{proof}
  This is the declared model definition of spring Constant In Newtons Per Meter.
\end{proof}
\begin{definition}[Figure Object]
  \label{def:physics:phyx-mini-0268:phyxminiproblems-problemphyxmini0268-figureobject}
  \lean{PhyXMiniProblems.ProblemPhyXMini0268.FigureObject}
  Objects visible in the supplied overhead-view image
\end{definition}
\begin{proof}
  This is the declared model definition of Figure Object.
\end{proof}
\begin{definition}[Figure Text Label]
  \label{def:physics:phyx-mini-0268:phyxminiproblems-problemphyxmini0268-figuretextlabel}
  \lean{PhyXMiniProblems.ProblemPhyXMini0268.FigureTextLabel}
  Text labels printed in the supplied image
\end{definition}
\begin{proof}
  This is the declared model definition of Figure Text Label.
\end{proof}
\begin{definition}[Figure Point]
  \label{def:physics:phyx-mini-0268:phyxminiproblems-problemphyxmini0268-figurepoint}
  \lean{PhyXMiniProblems.ProblemPhyXMini0268.FigurePoint}
  Distinguished points in the rod--spring diagram
\end{definition}
\begin{proof}
  This is the declared model definition of Figure Point.
\end{proof}
\begin{definition}[Motion Plane]
  \label{def:physics:phyx-mini-0268:phyxminiproblems-problemphyxmini0268-motionplane}
  \lean{PhyXMiniProblems.ProblemPhyXMini0268.MotionPlane}
  The plane in which the rod is free to move
\end{definition}
\begin{proof}
  This is the declared model definition of Motion Plane.
\end{proof}
\begin{definition}[Axis Orientation]
  \label{def:physics:phyx-mini-0268:phyxminiproblems-problemphyxmini0268-axisorientation}
  \lean{PhyXMiniProblems.ProblemPhyXMini0268.AxisOrientation}
  Orientation of the fixed rotation axis
\end{definition}
\begin{proof}
  This is the declared model definition of Axis Orientation.
\end{proof}
\begin{definition}[Rod Mass Model]
  \label{def:physics:phyx-mini-0268:phyxminiproblems-problemphyxmini0268-rodmassmodel}
  \lean{PhyXMiniProblems.ProblemPhyXMini0268.RodMassModel}
  The mass model for the rod
\end{definition}
\begin{proof}
  This is the declared model definition of Rod Mass Model.
\end{proof}
\begin{definition}[Rod Geometry]
  \label{def:physics:phyx-mini-0268:phyxminiproblems-problemphyxmini0268-rodgeometry}
  \lean{PhyXMiniProblems.ProblemPhyXMini0268.RodGeometry}
  The geometric idealization used for the rod's central moment of inertia
\end{definition}
\begin{proof}
  This is the declared model definition of Rod Geometry.
\end{proof}
\begin{definition}[Spring Force Model]
  \label{def:physics:phyx-mini-0268:phyxminiproblems-problemphyxmini0268-springforcemodel}
  \lean{PhyXMiniProblems.ProblemPhyXMini0268.SpringForceModel}
  The force model for the spring
\end{definition}
\begin{proof}
  This is the declared model definition of Spring Force Model.
\end{proof}
\begin{definition}[Rod Wall Alignment]
  \label{def:physics:phyx-mini-0268:phyxminiproblems-problemphyxmini0268-rodwallalignment}
  \lean{PhyXMiniProblems.ProblemPhyXMini0268.RodWallAlignment}
  Relative orientation of the rod and wall at equilibrium
\end{definition}
\begin{proof}
  This is the declared model definition of Rod Wall Alignment.
\end{proof}
\begin{definition}[Spring Rod Alignment]
  \label{def:physics:phyx-mini-0268:phyxminiproblems-problemphyxmini0268-springrodalignment}
  \lean{PhyXMiniProblems.ProblemPhyXMini0268.SpringRodAlignment}
  Relative orientation of the spring and rod at equilibrium
\end{definition}
\begin{proof}
  This is the declared model definition of Spring Rod Alignment.
\end{proof}
\begin{definition}[Release Condition]
  \label{def:physics:phyx-mini-0268:phyxminiproblems-problemphyxmini0268-releasecondition}
  \lean{PhyXMiniProblems.ProblemPhyXMini0268.ReleaseCondition}
  How the slightly displaced rod is started
\end{definition}
\begin{proof}
  This is the declared model definition of Release Condition.
\end{proof}
\begin{definition}[Supplied Rod Spring Figure]
  \label{def:physics:phyx-mini-0268:phyxminiproblems-problemphyxmini0268-suppliedrodspringfigure}
  \lean{PhyXMiniProblems.ProblemPhyXMini0268.SuppliedRodSpringFigure}
  \uses{def:physics:phyx-mini-0268:phyxminiproblems-problemphyxmini0268-figureobject, def:physics:phyx-mini-0268:phyxminiproblems-problemphyxmini0268-figuretextlabel, def:physics:phyx-mini-0268:phyxminiproblems-problemphyxmini0268-figurepoint}
  Incidence data and qualitative readouts transcribed from the primary image. The bitmap shows a displaced rod but contains no length scale, so depictedDisplacementAngle is a dimensionless schematic coordinate rather than measured numerical data
\end{definition}
\begin{proof}
  This is the declared model definition of Supplied Rod Spring Figure.
\end{proof}
\begin{definition}[Center Pivoted Rod Spring Setup]
  \label{def:physics:phyx-mini-0268:phyxminiproblems-problemphyxmini0268-centerpivotedrodspringsetup}
  \lean{PhyXMiniProblems.ProblemPhyXMini0268.CenterPivotedRodSpringSetup}
  \uses{def:physics:phyx-mini-0268:phyxminiproblems-problemphyxmini0268-massquantity, def:physics:phyx-mini-0268:phyxminiproblems-problemphyxmini0268-lengthquantity, def:physics:phyx-mini-0268:phyxminiproblems-problemphyxmini0268-signedlengthquantity, def:physics:phyx-mini-0268:phyxminiproblems-problemphyxmini0268-timequantity, def:physics:phyx-mini-0268:phyxminiproblems-problemphyxmini0268-springconstantquantity, def:physics:phyx-mini-0268:phyxminiproblems-problemphyxmini0268-signedforcequantity, def:physics:phyx-mini-0268:phyxminiproblems-problemphyxmini0268-momentofinertiaquantity, def:physics:phyx-mini-0268:phyxminiproblems-problemphyxmini0268-signedtorquequantity, def:physics:phyx-mini-0268:phyxminiproblems-problemphyxmini0268-restoringcoefficientquantity, def:physics:phyx-mini-0268:phyxminiproblems-problemphyxmini0268-angularaccelerationquantity, def:physics:phyx-mini-0268:phyxminiproblems-problemphyxmini0268-angularfrequencyquantity, def:physics:phyx-mini-0268:phyxminiproblems-problemphyxmini0268-motionplane, def:physics:phyx-mini-0268:phyxminiproblems-problemphyxmini0268-axisorientation, def:physics:phyx-mini-0268:phyxminiproblems-problemphyxmini0268-rodmassmodel, def:physics:phyx-mini-0268:phyxminiproblems-problemphyxmini0268-rodgeometry, def:physics:phyx-mini-0268:phyxminiproblems-problemphyxmini0268-springforcemodel, def:physics:phyx-mini-0268:phyxminiproblems-problemphyxmini0268-rodwallalignment, def:physics:phyx-mini-0268:phyxminiproblems-problemphyxmini0268-springrodalignment, def:physics:phyx-mini-0268:phyxminiproblems-problemphyxmini0268-releasecondition, def:physics:phyx-mini-0268:phyxminiproblems-problemphyxmini0268-suppliedrodspringfigure}
  Independent physical quantities and response fields of the experiment. The rod length is included even though the final closed form is length-independent. Neither the period nor the restoring coefficient is defined by the requested answer; governing laws below constrain them
\end{definition}
\begin{proof}
  This is the declared model definition of Center Pivoted Rod Spring Setup.
\end{proof}
\begin{definition}[Is In Small Angle Regime]
  \label{def:physics:phyx-mini-0268:phyxminiproblems-problemphyxmini0268-isinsmallangleregime}
  \lean{PhyXMiniProblems.ProblemPhyXMini0268.IsInSmallAngleRegime}
  \uses{def:physics:phyx-mini-0268:phyxminiproblems-problemphyxmini0268-centerpivotedrodspringsetup}
  The chosen neighborhood in which first-order angular laws are used
\end{definition}
\begin{proof}
  This is the declared model definition of Is In Small Angle Regime.
\end{proof}
\begin{definition}[Matches Problem And Figure Readouts]
  \label{def:physics:phyx-mini-0268:phyxminiproblems-problemphyxmini0268-matchesproblemandfigurereadouts}
  \lean{PhyXMiniProblems.ProblemPhyXMini0268.MatchesProblemAndFigureReadouts}
  \uses{def:physics:phyx-mini-0268:phyxminiproblems-problemphyxmini0268-massinkilograms, def:physics:phyx-mini-0268:phyxminiproblems-problemphyxmini0268-springconstantinnewtonspermeter, def:physics:phyx-mini-0268:phyxminiproblems-problemphyxmini0268-centerpivotedrodspringsetup}
  Numerical data, qualitative setup, and figure readouts supplied by the source. The rod length and pictured deflection have no numerical readout in the image. No field assigns the requested period or selects an answer choice
\end{definition}
\begin{proof}
  This is the declared model definition of Matches Problem And Figure Readouts.
\end{proof}
\begin{definition}[Has Physical Rod Spring Parameters]
  \label{def:physics:phyx-mini-0268:phyxminiproblems-problemphyxmini0268-hasphysicalrodspringparameters}
  \lean{PhyXMiniProblems.ProblemPhyXMini0268.HasPhysicalRodSpringParameters}
  \uses{def:physics:phyx-mini-0268:phyxminiproblems-problemphyxmini0268-massreadout, def:physics:phyx-mini-0268:phyxminiproblems-problemphyxmini0268-lengthreadout, def:physics:phyx-mini-0268:phyxminiproblems-problemphyxmini0268-timereadout, def:physics:phyx-mini-0268:phyxminiproblems-problemphyxmini0268-springconstantreadout, def:physics:phyx-mini-0268:phyxminiproblems-problemphyxmini0268-momentofinertiareadout, def:physics:phyx-mini-0268:phyxminiproblems-problemphyxmini0268-restoringcoefficientreadout, def:physics:phyx-mini-0268:phyxminiproblems-problemphyxmini0268-angularfrequencyreadout, def:physics:phyx-mini-0268:phyxminiproblems-problemphyxmini0268-centerpivotedrodspringsetup}
  Strict positivity and nondegeneracy assumptions. They make the rod and spring physical and support the scalar oscillator bridge without prescribing its period
\end{definition}
\begin{proof}
  This is the declared model definition of Has Physical Rod Spring Parameters.
\end{proof}
\begin{definition}[Satisfies Uniform Rod Inertia Law]
  \label{def:physics:phyx-mini-0268:phyxminiproblems-problemphyxmini0268-satisfiesuniformrodinertialaw}
  \lean{PhyXMiniProblems.ProblemPhyXMini0268.SatisfiesUniformRodInertiaLaw}
  \uses{def:physics:phyx-mini-0268:phyxminiproblems-problemphyxmini0268-massreadout, def:physics:phyx-mini-0268:phyxminiproblems-problemphyxmini0268-lengthreadout, def:physics:phyx-mini-0268:phyxminiproblems-problemphyxmini0268-momentofinertiareadout, def:physics:phyx-mini-0268:phyxminiproblems-problemphyxmini0268-centerpivotedrodspringsetup}
  For a uniform thin rod of length L rotating about a perpendicular axis through its center, I = m L² / 12. This local law predicate is needed because the available Physlib rigid-body rotational equation is informal and does not provide the uniform-rod formula
\end{definition}
\begin{proof}
  This is the declared model definition of Satisfies Uniform Rod Inertia Law.
\end{proof}
\begin{definition}[Satisfies Centered Endpoint Geometry]
  \label{def:physics:phyx-mini-0268:phyxminiproblems-problemphyxmini0268-satisfiescenteredendpointgeometry}
  \lean{PhyXMiniProblems.ProblemPhyXMini0268.SatisfiesCenteredEndpointGeometry}
  \uses{def:physics:phyx-mini-0268:phyxminiproblems-problemphyxmini0268-lengthreadout, def:physics:phyx-mini-0268:phyxminiproblems-problemphyxmini0268-centerpivotedrodspringsetup}
  Because the pivot is at the rod center and the spring is connected to one end, the spring's lever arm is L/2. This is figure/setup geometry, not a period formula
\end{definition}
\begin{proof}
  This is the declared model definition of Satisfies Centered Endpoint Geometry.
\end{proof}
\begin{definition}[Satisfies Linearized Rod Spring Laws]
  \label{def:physics:phyx-mini-0268:phyxminiproblems-problemphyxmini0268-satisfieslinearizedrodspringlaws}
  \lean{PhyXMiniProblems.ProblemPhyXMini0268.SatisfiesLinearizedRodSpringLaws}
  \uses{def:physics:phyx-mini-0268:phyxminiproblems-problemphyxmini0268-lengthreadout, def:physics:phyx-mini-0268:phyxminiproblems-problemphyxmini0268-signedlengthreadout, def:physics:phyx-mini-0268:phyxminiproblems-problemphyxmini0268-springconstantreadout, def:physics:phyx-mini-0268:phyxminiproblems-problemphyxmini0268-signedforcereadout, def:physics:phyx-mini-0268:phyxminiproblems-problemphyxmini0268-momentofinertiareadout, def:physics:phyx-mini-0268:phyxminiproblems-problemphyxmini0268-signedtorquereadout, def:physics:phyx-mini-0268:phyxminiproblems-problemphyxmini0268-restoringcoefficientreadout, def:physics:phyx-mini-0268:phyxminiproblems-problemphyxmini0268-angularaccelerationreadout, def:physics:phyx-mini-0268:phyxminiproblems-problemphyxmini0268-centerpivotedrodspringsetup, def:physics:phyx-mini-0268:phyxminiproblems-problemphyxmini0268-isinsmallangleregime}
  First-order mechanics for a signed angular displacement θ, whose positive direction is chosen to increase the spring length: the attached endpoint moves by r θ along the spring axis; Hooke's law gives force -k x; the spring torque is the force times the centered lever arm; the spring is the only source of torque in the horizontal plane; τ = -κ θ characterizes the effective restoring coefficient; rotational Newton's law is I θ'' = τ. These laws contain no period formula or numerical answer
\end{definition}
\begin{proof}
  This is the declared model definition of Satisfies Linearized Rod Spring Laws.
\end{proof}
\begin{definition}[scalar Rotational Oscillator]
  \label{def:physics:phyx-mini-0268:phyxminiproblems-problemphyxmini0268-scalarrotationaloscillator}
  \lean{PhyXMiniProblems.ProblemPhyXMini0268.scalarRotationalOscillator}
  \uses{def:physics:phyx-mini-0268:phyxminiproblems-problemphyxmini0268-momentofinertiareadout, def:physics:phyx-mini-0268:phyxminiproblems-problemphyxmini0268-restoringcoefficientreadout, def:physics:phyx-mini-0268:phyxminiproblems-problemphyxmini0268-centerpivotedrodspringsetup, def:physics:phyx-mini-0268:phyxminiproblems-problemphyxmini0268-hasphysicalrodspringparameters}
  The scalar equation I θ'' + κ θ = 0 has the same mathematical form as Physlib's positive harmonic oscillator. Its m slot is filled by the moment-of-inertia readout and its k slot by the restoring-torque coefficient readout; their ratio has the required inverse-time-squared dimension
\end{definition}
\begin{proof}
  This is the declared model definition of scalar Rotational Oscillator.
\end{proof}
\begin{definition}[Satisfies Ideal Small Oscillation Law]
  \label{def:physics:phyx-mini-0268:phyxminiproblems-problemphyxmini0268-satisfiesidealsmalloscillationlaw}
  \lean{PhyXMiniProblems.ProblemPhyXMini0268.SatisfiesIdealSmallOscillationLaw}
  \uses{def:physics:phyx-mini-0268:phyxminiproblems-problemphyxmini0268-timereadout, def:physics:phyx-mini-0268:phyxminiproblems-problemphyxmini0268-angularfrequencyreadout, def:physics:phyx-mini-0268:phyxminiproblems-problemphyxmini0268-centerpivotedrodspringsetup, def:physics:phyx-mini-0268:phyxminiproblems-problemphyxmini0268-hasphysicalrodspringparameters, def:physics:phyx-mini-0268:phyxminiproblems-problemphyxmini0268-scalarrotationaloscillator}
  The modeled angular frequency and period agree with Physlib's grounded ω = sqrt (k/m) and period = 2π/ω for the scalar rotational oscillator. This is a generic simple-harmonic-motion law; it does not expand I or κ using this problem's rod and spring parameters
\end{definition}
\begin{proof}
  This is the declared model definition of Satisfies Ideal Small Oscillation Law.
\end{proof}
\begin{lemma}[restoring Coefficient readout]
  \label{lem:physics:phyx-mini-0268:phyxminiproblems-problemphyxmini0268-restoringcoefficient-readout}
  \lean{PhyXMiniProblems.ProblemPhyXMini0268.restoringCoefficient_readout}
  \uses{def:physics:phyx-mini-0268:phyxminiproblems-problemphyxmini0268-lengthreadout, def:physics:phyx-mini-0268:phyxminiproblems-problemphyxmini0268-springconstantreadout, def:physics:phyx-mini-0268:phyxminiproblems-problemphyxmini0268-restoringcoefficientreadout, def:physics:phyx-mini-0268:phyxminiproblems-problemphyxmini0268-centerpivotedrodspringsetup, def:physics:phyx-mini-0268:phyxminiproblems-problemphyxmini0268-hasphysicalrodspringparameters, def:physics:phyx-mini-0268:phyxminiproblems-problemphyxmini0268-satisfieslinearizedrodspringlaws}
  The linearized spring geometry and Hooke law imply κ = k r². This is an intermediate consequence rather than a premise field
\end{lemma}
\begin{proof}
  The conclusion follows from the governing relations and figure readouts named in the statement of restoring Coefficient readout.
\end{proof}
\begin{definition}[Answer Choice]
  \label{def:physics:phyx-mini-0268:phyxminiproblems-problemphyxmini0268-answerchoice}
  \lean{PhyXMiniProblems.ProblemPhyXMini0268.AnswerChoice}
  The answer labels displayed with the exercise
\end{definition}
\begin{proof}
  This is the declared model definition of Answer Choice.
\end{proof}
\begin{definition}[displayed Period In Seconds]
  \label{def:physics:phyx-mini-0268:phyxminiproblems-problemphyxmini0268-displayedperiodinseconds}
  \lean{PhyXMiniProblems.ProblemPhyXMini0268.displayedPeriodInSeconds}
  \uses{def:physics:phyx-mini-0268:phyxminiproblems-problemphyxmini0268-answerchoice}
  Period values printed beside the answer choices, in seconds
\end{definition}
\begin{proof}
  This is the declared model definition of displayed Period In Seconds.
\end{proof}
\begin{definition}[recorded Dataset Answer]
  \label{def:physics:phyx-mini-0268:phyxminiproblems-problemphyxmini0268-recordeddatasetanswer}
  \lean{PhyXMiniProblems.ProblemPhyXMini0268.recordedDatasetAnswer}
  \uses{def:physics:phyx-mini-0268:phyxminiproblems-problemphyxmini0268-answerchoice}
  The answer label recorded in the source dataset; this is metadata, not a premise of either theorem below
\end{definition}
\begin{proof}
  This is the declared model definition of recorded Dataset Answer.
\end{proof}
\begin{definition}[Rounds To Four Decimal Places]
  \label{def:physics:phyx-mini-0268:phyxminiproblems-problemphyxmini0268-roundstofourdecimalplaces}
  \lean{PhyXMiniProblems.ProblemPhyXMini0268.RoundsToFourDecimalPlaces}
  A computed duration rounds to the displayed value at four decimal places
\end{definition}
\begin{proof}
  This is the declared model definition of Rounds To Four Decimal Places.
\end{proof}
\begin{definition}[Matches Answer Choice]
  \label{def:physics:phyx-mini-0268:phyxminiproblems-problemphyxmini0268-matchesanswerchoice}
  \lean{PhyXMiniProblems.ProblemPhyXMini0268.MatchesAnswerChoice}
  \uses{def:physics:phyx-mini-0268:phyxminiproblems-problemphyxmini0268-timeinseconds, def:physics:phyx-mini-0268:phyxminiproblems-problemphyxmini0268-centerpivotedrodspringsetup, def:physics:phyx-mini-0268:phyxminiproblems-problemphyxmini0268-answerchoice, def:physics:phyx-mini-0268:phyxminiproblems-problemphyxmini0268-displayedperiodinseconds, def:physics:phyx-mini-0268:phyxminiproblems-problemphyxmini0268-roundstofourdecimalplaces}
  A displayed choice agrees, to the printed precision, with the physical period's SI-second readout
\end{definition}
\begin{proof}
  This is the declared model definition of Matches Answer Choice.
\end{proof}
\begin{lemma}[small Oscillation Period formula]
  \label{lem:physics:phyx-mini-0268:phyxminiproblems-problemphyxmini0268-smalloscillationperiod-formula}
  \lean{PhyXMiniProblems.ProblemPhyXMini0268.smallOscillationPeriod_formula}
  \uses{def:physics:phyx-mini-0268:phyxminiproblems-problemphyxmini0268-massreadout, def:physics:phyx-mini-0268:phyxminiproblems-problemphyxmini0268-timereadout, def:physics:phyx-mini-0268:phyxminiproblems-problemphyxmini0268-springconstantreadout, def:physics:phyx-mini-0268:phyxminiproblems-problemphyxmini0268-centerpivotedrodspringsetup, def:physics:phyx-mini-0268:phyxminiproblems-problemphyxmini0268-hasphysicalrodspringparameters, def:physics:phyx-mini-0268:phyxminiproblems-problemphyxmini0268-satisfiesuniformrodinertialaw, def:physics:phyx-mini-0268:phyxminiproblems-problemphyxmini0268-satisfiescenteredendpointgeometry, def:physics:phyx-mini-0268:phyxminiproblems-problemphyxmini0268-satisfieslinearizedrodspringlaws, def:physics:phyx-mini-0268:phyxminiproblems-problemphyxmini0268-satisfiesidealsmalloscillationlaw}
  Uniform-rod inertia, centered endpoint geometry, the linearized spring torque, and the generic oscillator law imply T = 2π sqrt (m / (3k)). The rod length cancels between I = mL²/12 and κ = k(L/2)². No hypothesis contains this closed form
\end{lemma}
\begin{proof}
  The conclusion follows from the governing relations and figure readouts named in the statement of small Oscillation Period formula.
\end{proof}
% --- Archon declaration topology end ---
% --- Archon physics formalization source end ---
